%%% mode: latex
%%% TeX-master: t
%%% End:

% !!! 请用以下内容完全覆盖你的 body/conclusion.tex !!!

\chapter{总结与展望}
\label{chap:conclusion}

\section{本文工作总结}
雾天、霾等恶劣天气会导致图像对比度下降、颜色偏移以及目标边缘模糊，给户外的计算机视觉系统（如自动驾驶、无人机巡检及海事监控）带来了严峻挑战。现有的联合去雾与目标检测算法大多依赖于参数量庞大的深度卷积神经网络或复杂的 Transformer 架构，难以在算力受限的边缘设备上部署。此外，受限于真实雾天目标检测数据集（如 RTTS）的规模，复杂大模型在训练时极易陷入严重的过拟合陷阱。

针对上述问题，本文深入研究了雾天图像的退化机理与目标检测的特征提取机制，提出了一种基于状态空间模型（Mamba）与物理引导的轻量级联合目标检测算法（WF-DIDNet）。本文的主要研究工作和贡献总结如下：

\begin{enumerate}
    \item \textbf{验证了轻量级骨干网络在小样本雾天检测中的优越性。}
    本文打破了“大模型泛化能力必然更强”的思维惯性。通过在 RTTS 数据集上对 RT-DETR 框架进行基线测试，发现搭载 ResNet-101/50 的大模型在训练后期出现了明显的分类损失停滞与过拟合现象。相反，参数量较小的 ResNet-18 被迫学习更本质的抗雾语义特征，在 mAP50 指标上超越了大模型 1.5\%。这一工作不仅确立了本文的轻量化基准，也为后续恶劣天气小样本视觉任务的模型选择提供了重要参考。

    \item \textbf{设计了高分辨率 Mamba 去雾特征提取模块。}
    针对传统 CNN 编码器-解码器架构中下采样操作导致目标高频细节丢失的问题，本文引入了具有线性计算复杂度的状态空间模型（Mamba）。设计了无下采样操作的高分辨率去雾头（HighResMambaDehazeHead），利用选择性扫描机制高效捕捉全局上下文依赖，不仅成功输出了高保真的去雾重建图像，还在中间层截取了蕴含丰富大气透射率和浓度信息的物理特征图。

    \item \textbf{提出了基于动态通道适配的物理引导模块（PGM）。}
    为了解决传统级联算法中去雾特征利用不充分（即“图像级松耦合”）的问题，本文创新性地提出了特征级物理引导模块。该模块首先通过自动探测并动态重建卷积层，实现了去雾物理特征与检测主干语义特征（P3 层）的维度对齐；随后利用空间自适应全局平均池化与多层感知机（MLP），计算出通道级的注意力权重。通过重标定特征，网络能够精准“穿透”浓雾区域，显著提升了对模糊小目标的召回率。

    \item \textbf{制定了稳定高效的两阶段级联协同训练策略。}
    针对全新加入的去雾模块与预训练检测模块在端到端联合训练时容易导致梯度爆炸和权重破坏的问题，本文设计了先在 Haze4K 数据集上独立预训练去雾分支，再在 RTTS 检测数据集上微调的训练流水线。在微调阶段，严格冻结去雾网络的所有权重及批归一化（BN）层的统计量，保证了检测误差通过 PGM 回传时联合架构的高效、稳定收敛。
\end{enumerate}

综上所述，本文提出的 WF-DIDNet 在保证模型轻量化（参数量仅约 30M）的前提下，通过 Mamba 高效去雾与 PGM 物理引导的双重机制，显著提升了雾天场景下的目标检测精度，具有较高的学术创新性与工程实用价值。

\section{未来工作展望}
尽管本文提出的轻量级目标检测算法在雾天场景下取得了优异的性能，但受限于研究时间和计算资源，本研究仍存在一些局限性。未来的研究工作可以从以下几个方面展开深入探索：

\begin{enumerate}
    \item \textbf{复杂非均匀雾天与夜间场景的拓展}：
    本文的算法主要针对白天的均匀或半均匀雾天场景设计。在实际的无人机航拍或自动驾驶场景中，往往会遇到分布极不均匀的团雾，甚至伴随有城市路灯、车灯等强光源干扰的夜间带光晕雾天场景。未来的工作可以引入光照自适应评估模块或非均匀去雾物理先验，进一步提升模型在极端复杂光照及气象条件下的泛化能力。
    
    \item \textbf{底层硬件的部署与算子级优化}：
    虽然本文选用了轻量级的 ResNet-18 主干并控制了整体参数量，但 Mamba 架构作为一种相对较新的序列建模方法，其核心的“选择性扫描（Selective Scan）”算子在许多边缘计算设备（如无人机机载 NPU、NVIDIA Jetson 系列、FPGA 等）上的底层硬件加速支持尚不完善。未来的工作将侧重于模型量化（如 INT8 量化）、网络剪枝，以及对 SSM 核心算子进行 TensorRT 级别的定制开发，以实现算法在真实终端上的高帧率实时推理。
    
    \item \textbf{向多任务与时序多帧检测的延伸}：
    当前算法主要基于单幅图像进行静态目标检测。在实际监控与无人机巡航任务中，视频流数据包含了丰富的时序信息。未来可以尝试将当前框架与多目标跟踪（MOT）或实例分割任务相结合，并利用时序上下文（Temporal Context）来进一步弥补单帧图像在重雾遮挡下不可见特征的缺陷，构建更鲁棒的全天候视觉感知系统。
\end{enumerate}